\section{Discussion}

The experiments indicate that diffraction-domain loss alone can be insufficient
for judging reconstruction quality. A model may reduce FT L1 while producing an
incorrect object support or phase distribution. This is consistent with the
ill-posed nature of phase retrieval: multiple real-space configurations may
produce similar diffraction-domain errors unless additional priors are imposed.

\subsection{Why Far-field Fitting Can Be Misleading}
The far-field objective observes only the magnitude of the Fourier transform.
Consequently, it does not uniquely supervise the object-domain phase, support,
or spatial localization. The post-processing layer couples these variables, but
the loss is still applied after the FFT. A model can therefore exploit degrees
of freedom in amplitude and phase while still reducing the far-field residual.

\subsection{What the AutoPhaseNN Prior Provides}
The original AutoPhaseNN decoder appears to encode stronger spatial structure
through its architecture and post-processing assumptions. This prior may not be
a single explicit loss term. Instead, it can arise from the combination of
decoder topology, upsampling behavior, activation ranges, hard support masking,
and the oversampled diffraction setting. The practical effect is that the
network is biased toward compact object-domain solutions.

\subsection{Implication for Model Design}
The comparison suggests that replacing the backbone is not sufficient unless
the replacement preserves the physical bias of the original pipeline. A useful
future model should either inherit a hard spatial prior from its architecture or
enforce the prior through a simple post-processing constraint. Soft auxiliary
losses may help optimization, but they should not be presented as the central
solution unless they demonstrably improve real-space amplitude, phase, and
support at the same time.

\todo{Add concise analysis of the square/central support prior and oversampling
assumption, but keep the discussion focused on what is needed for the proposed
model design.}
